
\documentclass[tikz,border=5pt]{standalone}
\usepackage{tikz}



\begin{document}




\def\DX{4}
\def\DY{2}
\def\DZ{4}

\definecolor{c0}{HTML}{1F77B4} % red
\definecolor{c1}{HTML}{2CA02C} % green
\definecolor{c2}{HTML}{D62728} % blue
\definecolor{c3}{HTML}{9467BD} % purple
\definecolor{c4}{HTML}{17BECF} % cyan
\definecolor{c5}{HTML}{FF7F0E} % orange
\definecolor{c6}{HTML}{8C564B}





\newcommand{\colorbyy}[1]{%
\ifcase#1 c0\or c1\or c2\or c3\or c4\or c5\fi
}

\newcommand{\colorbyyRev}[1]{%
\ifcase#1 c5\or c4\or c3\or c2\or c1\or c0\fi
}

\newcommand{\colorbyyRevY}[1]{%
\ifcase#1 c2\or c1\or c0\fi
}


\newcommand{\rotshapeforward}[3]{
  % input: #1=x, #2=y, #3=z
  % output: rotated indices in \outx, \outy, \outz

  \edef\outx{#1}% X unchanged
  \edef\outy{#3}% new Y = old Z
  \edef\outz{#2}% new Z = old Y
}




\newcommand{\cubeview}[4]{ % #1: fill color, #2,#3,#4: x,y,z coords
    \begin{scope}[shift={(#2,#3,#4)}]
        \draw[fill=#1!50, draw=black] (0,0,1) -- (1,0,1) -- (1,1,1) -- (0,1,1) -- cycle;
        \draw[fill=#1!70, draw=black] (1,0,0) -- (1,1,0) -- (1,1,1) -- (1,0,1) -- cycle;
        \draw[fill=#1!70, draw=black] (0,0,0) -- (1,0,0) -- (1,0,1) -- (0,0,1) -- cycle;
    \end{scope}
}



\begin{tikzpicture}[
    x={(1cm,0cm)},
    y={(0.6cm,0.35cm)},
    z={(0cm,1cm)},
    line join=round,
    line width=0.4pt
]

\rotshapeforward{\DX}{\DY}{\DZ}


\pgfmathtruncatemacro{\outyminusone}{\outy-1}

\foreach \z in {0,1,...,\outz}{
    \foreach \x in {0,1,...,\outx}{
        \foreach \y in {\outy,\outyminusone,...,0}{
            \cubeview{\colorbyy{\x}}{\x}{\y}{\z}
        }
    }
}


\end{tikzpicture}




\end{document}